\color{unidarkgreen}{\subsection{Блокировка при неисправностях в цепях напряжения (БНН) 2СШ}}
\color{black}

\begin{enumerate}[label=\arabic{section}.\arabic{subsection}.\arabic*,  align=left, labelsep=4pt, leftmargin=0pt, itemindent=57pt]

\item
Для алгоритмов РЗА, которые при неисправностях в цепях напряжения ТН 2СШ  могут излишне или
ложно сработать и повлиять на режим работы первичной сети, используется специальная блокировка при
неисправностях в цепях напряжения (БНН), которая вводится программным ключом \textbf{«Контр. ЦН 2СШ»}.

\item
В качестве основного алгоритма БНН используется критерий пофазного сравнения напряжений основной
вторичной обмотки, собранной по схеме «звезда» (ВО-1), и дополнительной вторичной обмоткой, собранной
по схеме «разомкнутый треугольник» (ВО-2). Между замерами цепей «звезды» и приведенных к ним замерами
«разомкнутого треугольника» производится расчёт векторных невязок по следующим формулам:

\begin{equation}\label{eq:BNN_a}
U_{\textsl{БНН ф.A}} = \left| \vec U_{A\textsl{(ВО-1)}} - \frac{\vec U_{a\textsl{(ВО-2)}}}{\sqrt{3}} \right|
\end{equation}

\begin{equation}\label{eq:BNN_b}
U_{\textsl{БНН ф.В}} = \left| \vec U_{B\textsl{(ВО-1)}} - \frac{\vec U_{b\textsl{(ВО-2)}}}{\sqrt{3}} \right|
\end{equation}

\begin{equation}\label{eq:BNN_c}
U_{\textsl{БНН ф.С}} = \left| \vec U_{C\textsl{(ВО-1)}} - \frac{\vec U_{c\textsl{(ВО-2)}}}{\sqrt{3}} \right|
\end{equation}

\item БНН по критерию пофазного сравнения напряжений вводится в работу программным ключом \textbf{«БНН-ф
2СШ»}. Неисправность в цепях напряжения определяется, если величина векторной невязки по одной из фаз 
превышает уставку \textbf{«Uср БНН-ф 2СШ»}. 

\item В зависимости от схемы подключения дополнительной вторичной обмотки с помощью программных 
переключателей \textbf{«Ua ВО-2»}, \textbf{«Ub ВО-2»} и \textbf{«Uc ВО-2»} можно задать, какой вывод «разомкнутого треугольника», 
соответствует фазе «звезды». Для предотвращения ложного срабатывания БНН-ф при однофазном КЗ внутри 
контура заземления подстанции (в случае неправильного выполнения цепей заземления) предусмотрена 
блокировка по току нулевой последовательности, которая вводится программным ключом 
\textbf{«Контр. 3I0 БНН-ф 2СШ»}. Уставка блокировки по току нулевой последовательности не регулируется и равна Iном.

\item 
В устройстве предусмотрен алгоритм БНН, повторяющий принцип построения КРБ-12М, который
реагирует на напряжение небаланса $U_{\textsl{БНН}}$, В, между основной и дополнительной вторичными обмотками. 

\item 
Расчет напряжения небаланса зависит от схемы подключения цепей «разомкнутого треугольника». В таблице
\ref{bnn:tbl} приведены уставки для БНН по сравнению напряжений «звезды» и «разомкнутого треугольника».

%%%%%%%%%%%%%ТАБЛИЦА %%%%%%%%%%%%%%%%%%%
\small
\setlength{\extrarowheight}{0.05cm} 
\begin{longtable}
{|>{\centering\arraybackslash}m{6cm}|>{\centering\arraybackslash}m{4cm}|>{\centering\arraybackslash}m{2cm}|>{\centering\arraybackslash}m{4cm}|}
\caption{Перечень уставок алгоритма БНН по принципу КРБ-12М\hfill\vspace{-0.5\baselineskip}}\label{bnn:tbl}\\ 
\hhline{|-|-|-|-|}
\rowcolor{gray!30}
Наименование уставки  & Диапазон   & Шаг & Значение по умолчанию  \\ 
\hhline{|-|-|-|-|}
\endfirsthead
\caption*{\hspace{3pt}\emph{Продолжение таблицы \ref{bnn:tbl}\hfill\vspace{-0.5\baselineskip}}} \\ % сделано по ГОСТ 2.105 п.6.8.7
\hhline{|-|-|-|-|}
\rowcolor{gray!30}
Наименование уставки  & Диапазон   & Шаг & Значение по умолчанию  \\ 
\hhline{|-|-|-|-|}
\endhead
\hhline{|-|-|-|-|}
\endfoot
\endlastfoot
\centering Схема подключения цепей разомкнутого треугольника & \centering Uни/Uик \\ Uнф/Uфк \\ Uни/Uиф/Uфк \\ Uнк & \centering - & \centering \arraybackslash Uни/Uиф/Uфк \\  
\hline
\centering Особая фаза ТН  & \centering Выведено & \centering - & \centering \arraybackslash А \\  
\hline
\centering Направление векторов звезды и треугольника  & \centering Совпадает/не совпадает  & \centering - & \centering\arraybackslash Совпадает \\  
\hline
\end{longtable}
\setlength{\extrarowheight}{0.0cm} % возврат параметра высота ячеек в таблице
\normalsize    

\item
Схемы соединения обмоток, формулы для расчета $U_{\textsl{ср} ИО БНН 2СШ}$, векторные диаграммы приведены в
\textcolor{red}{приложении \ref{app:bnn}}.

\item
Алгоритм, повторяющий принцип построения КРБ-12М, вводится в работу программным ключом \textbf{«ИО БНН 2СШ»}.
Порог срабатывания по небалансу напряжений задается уставкой \textbf{«Uср ИО БНН 2СШ»}.

\item 
Также предусмотрен алгоритм БНН, основанный на комбинированном принципе, контролирующий
напряжения и токи. Алгоритм вводится программным ключом \textbf{«Комб. БНН 2СШ»}. Для выявления
несимметричных повреждений в цепях напряжения (обрыв или короткое замыкание одной или двух фаз)
контролируются следующие условия:
\begin{itemize}
\item   превышение напряжения обратной последовательности заданного значения (более 0,085$\cdot U_{\textsl{ном}}$);
\item   отсутствие тока обратной последовательности. БНН блокируется при превышении отношения тока обратной последовательности к току прямой последовательности заданной уставки \textbf{«I2/I1 комб. БНН 2СШ»} -- это необходимо для исключения излишнего срабатывания БНН при несимметричных коротких замыканиях в силовой цепи. Рекомендуемым значением коэффициента является 0,2;
\item   контроль отношения напряжения обратной последовательности к напряжению прямой последовательности. БНН срабатывает только при превышении отношения заданной уставки \textbf{«U2/U1 комб. БНН 2СШ»}. Рекомендуемым значением коэффициента является 0,2;
\item   отсутствие броска тока прямой последовательности. БНН блокируется при превышении приращения тока прямой последовательности заданной уставки \textbf{«dI1 комб. БНН 2СШ»}.
\end{itemize}

\item 
Обрыв трех фаз как во вторичной цепи, определяется при одновременном выполнении следующих условий:
\begin{itemize}
    \item   скачкообразное снижение напряжения прямой последовательности более чем на 0,2$\cdot$Uном;
    \item   отсутствие броска тока прямой последовательности, превышающего заданную величину, задаваемую уставкой \textbf{«dI1 комб. БНН 2СШ»}.
\end{itemize}

\item
При этом переход несимметричного повреждения в симметричное не приводит к возврату БНН, а также
обеспечивается отсутствие ложного срабатывания при КЗ в силовой цепи.

\item
Сигнал обрыва трех фаз подхватывается энергонезависимым триггером при снижении напряжени
прямой последовательности ниже уставки \textbf{«U1 мин БНН 2СШ»}. Возврат цепи происходит после повторного
повышения напряжения прямой последовательности более уставки \textbf{«U1 мин БНН 2СШ»} (рекомендуемое значение
уставки 0,1$\cdot$Uном).

\item
Для контроля одновременного исчезновения трех фазных напряжений используются три реле
минимального напряжения в фазах А, В, С. Отсутствие напряжения определяется при снижении всех фазных
напряжений меньше уставки \textbf{«Uф мин БНН 2СШ»} в течение 5 секунд и наличии тока в защищаемой стороны.
Отсутствие напряжение действует на формирование сигнала обрыва трех фаз.

\item
Логика БНН включает в себя контроль обрыва нуля схемы «звезда». Для выявления обрыва нулевого
провода предусмотрено создание искусственной несимметрии путём подключения догрузочного резистора к
входу фазы В в измерительных цепях напряжения. Обрыв нулевого провода определяется при превышении
напряжения нулевой последовательности 0,085$\cdot$Uном и отсутствии появления тока нулевой
последовательности, контролируемого уставкой \textbf{«3I0/I1 БНН 2СШ»}. Данный алгоритм также реагирует на обрыв
одной или двух фаз цепей напряжения. Ввод алгоритма осуществляется программным ключом
\textbf{«Контр. обрыв НП 2СШ»}.

\item 
Предусмотрена команда управления \textit{«Фикс. Неиспр. ЦН 2СШ»} для принудительного ввода состояния неисправности измерительных цепей напряжения.

\item
Для контроля состояния автомата ТН предусмотрен вход ФСУ \textit{«Автомат ТН 2СШ»}. Сигнал может
приниматься как от нормально открытых, так и от нормально закрытых блок-контактов. Выбор типа блок-
контактов осуществляется программным переключателем \textbf{«Конт. Автомат ТН 2СШ»}. Неудачная попытка
восстановления напряжения при включении автомата ТН на КЗ не приводит к возврату БНН.

\item
Действие БНН осуществляется без выдержки времени на блокировку зависимых алгоритмов РЗА, а с
выдержкой времени \textbf{«Тсигн неиспр. ЦН 2СШ»} происходит действие на сигнализацию.

\item
Оперативный ввод/вывод БНН 2СШ производится с помощью БУ (блоков управления функциями) \textit{«Вывод БНН 2СШ}.

\begin{figure}[H]
\centering
\includegraphics[width=0.85\textwidth,height=0.85\textheight,keepaspectratio]{\fbpath/661.1472 БНН 2СШ/img/BNN.pdf}
\captionof{figure}{Функциональная схема БНН}\label{pict_bnn:BNN}
\end{figure}

\small
\begin{longtable}{|>{\centering\arraybackslash}m{5.3cm}|>{\centering\arraybackslash}m{3.3cm}|>{\centering\arraybackslash}m{4.2cm}|>{\centering\arraybackslash}m{1.8cm}|>{\centering\arraybackslash}m{1cm}|}
\caption{Перечень уставочных параметров БНН\hfill\vspace{-0.5\baselineskip}}\label{bnn:tbl2}\\ 
\hline
\rowcolor{gray!30}
Параметр & Условное обозначение на схеме & Значение/ Диапазон & Единица измерения & Шаг \\ 
\hline
\endfirsthead
\caption*{\hspace{3pt}\emph{Продолжение таблицы \ref{bnn:tbl2}\hfill\vspace{-0.5\baselineskip}}} \\ % сделано по ГОСТ 2.105 п.6.8.7
\hline
\rowcolor{gray!30}
Параметр & Условное обозначение на схеме & Значение/ Диапазон & Единица измерения & Шаг \\ 
%\hline
\endhead

%===m>БНН|-
%===m
\end{longtable}
\normalsize

\end{enumerate}